\BeleskeID{04-s016}
% element: 04-s016:e01 — Odnos osnova p=q^k i transformacija svake cifre
% element: 04-s016:e02 — Razvijanje CVp pomoću p i zamena stepenima q
% element: 04-s016:e03 — Skupovi cifara p i q; razrada jedne cifre preko k manjih cifara
% element: 04-s016:e04 — Završni izraz CVq sa dvostrukim indeksima u izvornom zapisu
% element: 04-s016:e05 — Obrnuti prelazak grupisanjem po k cifara
% element: 04-s016:d01

Odnos $p=q^k$ omogućava da jednu cifru veće osnove zamenimo grupom tačno $k$ cifara manje osnove. Kada razvijanje jedne cifre pomnožimo težinom njenog mesta $p^i=q^{ki}$, dobijamo
\[c_i p^i=(b_{k-1}q^{k-1}+b_{k-2}q^{k-2}+\cdots+b_0)q^{ki}.\]
Njene težine zato postaju $q^{ki+k-1},q^{ki+k-2},\ldots,q^{ki}$. Sledeća cifra zauzima narednu grupu težina, bez preklapanja grupa. Ne menjamo brojnu vrednost, već isti zbir pišemo sa manjom osnovom.

U svakoj grupi zadržavamo $k$ mesta, uključujući eventualne vodeće nule. Obrnuto, grupa od $k$ cifara manje osnove može imati vrednost najviše $q^k-1=p-1$, pa uvek odgovara jednoj dozvoljenoj cifri veće osnove.
